\chapter*{Abstract}
\addcontentsline{toc}{chapter}{Abstract}

Graph neural networks are increasingly used for clinical prediction from
electronic health records, but they are usually explained after the fact, by
post-hoc attribution methods whose output cannot be checked against the way the
model actually decided. This project asks whether a self-explainable
architecture is a workable alternative in that setting. We build one
heterogeneous graph per emergency-department visit from MIMIC-IV-ED, with
vital signs, medications, symptoms and chief complaints as typed nodes, and
predict the primary diagnosis across 30 merged disease classes. Predictions are
made by a prototype layer that compares each patient graph to 90 learned
archetypes, which are projected onto real training subgraphs. The model is
compared against a knowledge-graph GNN (GraphCare), a prototype-free ablation of
the same encoder, and a tabular gradient-boosting baseline, all on one canonical
subject-aware split.

Three results stand out. The tabular baseline is the strongest predictor on
every metric (65.4\% accuracy, 0.607 macro-F1) against 56.3\% and 0.508 for the
prototype model, so the graph structure used here adds interpretability rather
than accuracy. The prototype layer itself is close to free: removing it changes
accuracy by 0.2 percentage points. And \caveat{on a 50-graph sample of the test fold}, the three post-hoc
explainers agree on the most important node in only
32\% of the cases, while the prototype evidence matches the prediction in 39 of
those 50 and its distance separates correct from incorrect predictions, which
argues for building the explanation into the model.

\tbd{A third comparison method (Method~C) is planned but not yet implemented.
Once it is, add one sentence here describing what it contributes and how it
scores, and update the numbers above if it changes the ordering.}
